

The exact minimax rate for recovering the conductivity \(f\) in this Darcy problem is **not currently established** in the literature. The best rigorously proved upper bound is
\[
n^{-\frac{2(\alpha-1)}{2\alpha+2+d}}
\]
(up to logarithmic factors), and this is the conjectured minimax rate. However, no matching minimax lower bound for \(f\) is available in the cited references; only the prediction problem for \(u_f\) is known to be minimax optimal.

Sources:
- [iopscience.iop.org](https://google.iopscience.iop.org/article/10.1088/1361-6420/ab7d2a)
- [exa.ai](https://exa.ai/library/publication/rj7jk1yjt53)
- [arxiv.org](https://arxiv.org/html/2601.09007)
- [arxiv.org](https://export.arxiv.org/pdf/1809.08818v3.pdf)
- [arxiv.org](https://arxiv.org/html/2409.03417)

Below is a compilable LaTeX document giving the rigorous upper bound and explaining why the minimax statement is presently open.

```latex
\documentclass[11pt]{article}

\usepackage{amsmath,amssymb,amsthm,mathtools,hyperref}

\theoremstyle{plain}
\newtheorem{theorem}{Theorem}
\newtheorem{proposition}[theorem]{Proposition}
\newtheorem{lemma}[theorem]{Lemma}
\newtheorem{corollary}[theorem]{Corollary}

\theoremstyle{definition}
\newtheorem{remark}[theorem]{Remark}

\newcommand{\R}{\mathbb R}
\newcommand{\E}{\mathbb E}
\newcommand{\1}{\mathbf 1}
\newcommand{\F}{\mathcal F}
\newcommand{\D}{D}
\DeclareMathOperator{\divg}{div}

\title{On the minimax rate for a Darcy conductivity inverse problem in Gaussian white noise}
\author{}
\date{}

\begin{document}

\maketitle

\begin{abstract}
We consider the inverse problem of estimating the conductivity \(f\) in
\[
\divg(f\nabla u_f)=1 \quad\text{in }D,\qquad u_f|_{\partial D}=0,
\]
from Gaussian white-noise observations of \(u_f\). We show that the best currently proved statistical upper bound for the \(L^2\)-risk is
\[
n^{-\frac{2(\alpha-1)}{2\alpha+2+d}}
\]
(up to logarithmic factors), under the natural smoothness assumption \(f\in C^\alpha\), \(\alpha>1\), and for conductivities sufficiently close to \(1\). We also explain that a matching minimax lower bound for \(f\) is not presently available in the literature; consequently the exact minimax rate is open. The conjectured minimax rate is the above upper bound.
\end{abstract}

\section{Problem and main conclusion}

Let \(D\subset \R^d\), \(d\ge 2\), be the Euclidean unit ball. For \(f>0\), let \(u_f\) be the unique weak solution of
\[
\divg(f\nabla u_f)=1 \quad\text{in }D,\qquad u_f|_{\partial D}=0.
\]
We observe
\[
Y_n = u_f + n^{-1/2}\mathbb W,
\]
where \(\mathbb W\) is Gaussian white noise on \(L^2(D)\). For \(\alpha>1\) and fixed small \(r\in(0,1)\), consider
\[
\F(\alpha,r)
=
\left\{
h\in C^\alpha(D): \|h-\1\|_{C^\alpha(D)}\le r
\right\}.
\]
The minimax risk is
\[
R_n^*(\alpha,r)
=
\inf_{\hat f_n}
\sup_{f_0\in \F(\alpha,r)}
\E_{f_0}\bigl[\|\hat f_n-f_0\|_{L^2(D)}^2\bigr].
\]

The central rigorous conclusion is the following.

\begin{theorem}[State of the art]
\label{thm:state}
Let \(\alpha>1\), and let \(r>0\) be sufficiently small. Then there exists an estimator \(\hat f_n\) such that, as \(n\to\infty\),
\[
\sup_{f_0\in\F(\alpha,r)}
\E_{f_0}\bigl[\|\hat f_n-f_0\|_{L^2(D)}^2\bigr]
\le
C\, n^{-\frac{2(\alpha-1)}{2\alpha+2+d}}
\]
(up to logarithmic factors if the smoothness is unknown). Here \(C\) depends only on \(d,\alpha,r\).

At present, no matching minimax lower bound for estimating \(f\) in \(L^2(D)\) is known in general for this model. The prediction problem,
\[
\inf_{\tilde u_n}
\sup_{f_0\in\F(\alpha,r)}
\E_{f_0}\bigl[\|\tilde u_n-u_{f_0}\|_{L^2(D)}^2\bigr],
\]
is known to have minimax rate
\[
n^{-\frac{2(\alpha+1)}{2\alpha+2+d}}.
\]
Thus the conjectured minimax rate for estimating \(f\) is
\[
R_n^*(\alpha,r)
\asymp
n^{-\frac{2(\alpha-1)}{2\alpha+2+d}}
\]
(up to logarithmic factors), but this is not presently a theorem in full generality.
\end{theorem}

The remainder gives a proof of the upper bound and discusses the lower-bound status.

\section{Preliminaries}

We collect the analytic facts needed for the proof.

\begin{lemma}[Uniform ellipticity and well-posedness]
\label{lem:wellposed}
If \(f\in \F(\alpha,r)\) with \(r<1\), then
\[
1-r \le f(x)\le 1+r\qquad x\in D.
\]
Hence for every \(f\in\F(\alpha,r)\) the boundary value problem
\[
\divg(f\nabla u_f)=1\quad\text{in }D,\qquad u_f|_{\partial D}=0
\]
has a unique weak solution \(u_f\in H^1_0(D)\), and the solution depends continuously on \(f\).
\end{lemma}

\begin{proof}
This is the standard Lax--Milgram theory for uniformly elliptic divergence-form equations. The uniform ellipticity follows from \(\|f-\1\|_{L^\infty}\le r<1\).
\end{proof}

\begin{lemma}[Elliptic regularity]
\label{lem:reg}
Let \(\alpha>1\). There exists \(M=M(d,\alpha,r)\) such that
\[
\sup_{f\in\F(\alpha,r)}\|u_f\|_{C^{\alpha+1}(\overline D)}
\le M.
\]
In particular, the forward image
\[
\mathcal U(\alpha,r):=\{u_f:f\in\F(\alpha,r)\}
\]
is contained in a fixed Hölder--Zygmund ball \(C^{\alpha+1}(\overline D)\).
\end{lemma}

\begin{proof}
Since \(f\in C^\alpha\) and \(f\ge 1-r>0\), the equation can be written locally in non-divergence form as
\[
f\Delta u_f + \nabla f\cdot \nabla u_f = 1.
\]
Because \(\alpha>1\), \(\nabla f\in C^{\alpha-1}\). Standard Schauder estimates for uniformly elliptic equations with Hölder continuous coefficients give \(u_f\in C^{\alpha+1}\), uniformly over the class \(\F(\alpha,r)\). See, for example, the classical theory in Ladyzhenskaya--Ural'tseva or Gilbarg--Trudinger. The boundary is smooth, and the boundary data are zero, so the estimate holds up to the boundary.
\end{proof}

\begin{lemma}[One-measurement stability for Darcy]
\label{lem:stability}
There exist constants \(r_0>0\) and \(C=C(d,\alpha,r_0)\) such that, for every \(r\le r_0\) and every \(f_1,f_2\in\F(\alpha,r)\),
\[
\|f_1-f_2\|_{L^2(D)}
\le
C\|u_{f_1}-u_{f_2}\|_{H^2(D)}.
\]
\end{lemma}

\begin{proof}
This is a standard stability estimate for the Darcy inverse problem with one interior measurement and known source. It follows by subtracting the two equations,
\[
\divg(f_1\nabla u_{f_1})=1,
\qquad
\divg(f_2\nabla u_{f_2})=1,
\]
and using that, for conductivities sufficiently close to \(1\), the linearized map
\[
h\mapsto v,\qquad
\Delta v + \divg(h\nabla u_{\1})=0,
\]
has a stable inverse from \(H^2\) to \(L^2\). A complete proof is given in the PDE-constrained inverse problem literature; see for instance the stability verification for the divergence-form Darcy problem in
\href{https://export.arxiv.org/pdf/1809.08818v3.pdf}{arxiv.org}
and related references therein.
\end{proof}

\section{Estimation of \(u_f\) in \(H^2\)}

The following lemma is the nonparametric regression ingredient.

\begin{lemma}
\label{lem:u-h2-rate}
Let \(\beta>2\). Suppose \(u\) belongs to a bounded subset of \(C^\beta(D)\). In the white-noise model
\[
Y_n = u + n^{-1/2}\mathbb W,
\]
there exists an estimator \(\tilde u_n\) such that
\[
\sup_{u\in C^\beta,\ \|u\|_{C^\beta}\le M}
\E\bigl[\|\tilde u_n-u\|_{H^2(D)}^2\bigr]
\le
C n^{-\frac{2(\beta-2)}{2\beta+d}}.
\]
If \(\beta\) is unknown, the same rate holds up to logarithmic factors.
\end{lemma}

\begin{proof}
Use a compactly supported wavelet basis adapted to the smooth domain \(D\). Let \(P_J\) denote the \(L^2\)-orthogonal projection onto wavelet levels \(j\le J\). Define
\[
\tilde u_n=P_JY_n.
\]
Write
\[
\tilde u_n-u = P_J(u+n^{-1/2}\mathbb W)-u
= -(I-P_J)u + n^{-1/2}P_J\mathbb W.
\]

For \(u\in C^\beta=B_{\infty,\infty}^\beta\), the wavelet coefficients satisfy
\[
|\langle u,\psi_{j,k}\rangle|
\lesssim
2^{-j(\beta+d/2)}.
\]
Using the wavelet characterization of \(H^2\),
\[
\|v\|_{H^2}^2
\asymp
\sum_{j,k}2^{4j}|\langle v,\psi_{j,k}\rangle|^2,
\]
we obtain the bias bound
\[
\|(I-P_J)u\|_{H^2}^2
\lesssim
\sum_{j>J}2^{4j}2^{jd}2^{-2j(\beta+d/2)}
=
\sum_{j>J}2^{-2j(\beta-2)}
\lesssim
2^{-2J(\beta-2)}.
\]

For the stochastic term,
\[
\E\|P_J\mathbb W\|_{H^2}^2
\asymp
\sum_{j\le J}2^{4j}2^{jd}
\lesssim
2^{J(d+4)}.
\]
Therefore
\[
\E\|\tilde u_n-u\|_{H^2}^2
\lesssim
2^{-2J(\beta-2)} + n^{-1}2^{J(d+4)}.
\]
Choosing
\[
2^{J_n}\asymp n^{1/(2\beta+d)}
\]
balances the two terms and gives
\[
\E\|\tilde u_n-u\|_{H^2}^2
\lesssim
n^{-\frac{2(\beta-2)}{2\beta+d}}.
\]
If \(\beta\) is unknown, a standard Lepski or thresholding procedure yields the same rate up to logarithmic factors.
\end{proof}

Applying Lemma~\ref{lem:u-h2-rate} with \(\beta=\alpha+1\) gives the following.

\begin{corollary}
\label{cor:u-h2}
There exists an estimator \(\tilde u_n\) such that
\[
\sup_{f_0\in\F(\alpha,r)}
\E_{f_0}\bigl[\|\tilde u_n-u_{f_0}\|_{H^2(D)}^2\bigr]
\le
C n^{-\frac{2(\alpha-1)}{2\alpha+2+d}},
\]
up to logarithmic factors if \(\alpha\) is unknown.
\end{corollary}

\section{Upper bound for estimating \(f\)}

We now construct an estimator of \(f\) from the estimator \(\tilde u_n\).

Let
\[
G:\F(\alpha,r)\to H^2(D),\qquad G(f)=u_f.
\]
By Lemma~\ref{lem:reg} and standard compactness arguments, \(G(\F(\alpha,r))\) is relatively compact in \(H^2(D)\). Define \(\hat f_n\in\F(\alpha,r)\) as an approximate minimizer of
\[
f\mapsto \|u_f-\tilde u_n\|_{H^2(D)}.
\]
More precisely, choose \(\hat f_n\) such that
\[
\|u_{\hat f_n}-\tilde u_n\|_{H^2}
\le
\inf_{f\in\F(\alpha,r)}\|u_f-\tilde u_n\|_{H^2}
+ n^{-1}.
\]

\begin{theorem}[Achievable rate]
\label{thm:upper}
Let \(\alpha>1\) and let \(r>0\) be sufficiently small. The estimator \(\hat f_n\) satisfies
\[
\sup_{f_0\in\F(\alpha,r)}
\E_{f_0}\bigl[\|\hat f_n-f_0\|_{L^2(D)}^2\bigr]
\le
C n^{-\frac{2(\alpha-1)}{2\alpha+2+d}},
\]
up to logarithmic factors.
\end{theorem}

\begin{proof}
Fix \(f_0\in\F(\alpha,r)\). By the approximate minimizing property,
\[
\|u_{\hat f_n}-\tilde u_n\|_{H^2}
\le
\|u_{f_0}-\tilde u_n\|_{H^2}+n^{-1}.
\]
Therefore
\[
\|u_{\hat f_n}-u_{f_0}\|_{H^2}
\le
\|u_{\hat f_n}-\tilde u_n\|_{H^2}
+
\|\tilde u_n-u_{f_0}\|_{H^2}
\le
2\|\tilde u_n-u_{f_0}\|_{H^2}+n^{-1}.
\]
Using the stability estimate Lemma~\ref{lem:stability},
\[
\|\hat f_n-f_0\|_{L^2}
\le
C\|u_{\hat f_n}-u_{f_0}\|_{H^2}
\le
C\left(2\|\tilde u_n-u_{f_0}\|_{H^2}+n^{-1}\right).
\]
Squaring and taking expectation,
\[
\E_{f_0}\|\hat f_n-f_0\|_{L^2}^2
\le
C\E_{f_0}\|\tilde u_n-u_{f_0}\|_{H^2}^2
+
C n^{-2}.
\]
Taking the supremum over \(f_0\in\F(\alpha,r)\) and applying Corollary~\ref{cor:u-h2} yields the result.
\end{proof}

\section{Lower bounds and minimax status}

For the prediction problem, the minimax rate is known. In the present notation,
\[
\inf_{\tilde u_n}
\sup_{f_0\in\F(\alpha,r)}
\E_{f_0}\bigl[\|\tilde u_n-u_{f_0}\|_{L^2(D)}^2\bigr]
\asymp
n^{-\frac{2(\alpha+1)}{2\alpha+2+d}}.
\]
This follows from the PDE-constrained regression theory for divergence-form equations; see
\href{https://exa.ai/library/publication/rj7jk1yjt53}{exa.ai}
and
\href{https://export.arxiv.org/pdf/1809.08818v3.pdf}{arxiv.org}.

For the recovery of \(f\) itself, however, a matching lower bound is not available in the cited literature. In particular, the Bayesian inverse-problem analysis in
\href{https://google.iopscience.iop.org/article/10.1088/1361-6420/ab7d2a}{iopscience.iop.org}
explicitly notes that no matching lower bounds are known for the conductivity recovery rates. Likewise, the later computational/statistical analysis in
\href{https://arxiv.org/html/2601.09007}{arxiv.org}
obtains the same upper rate and notes that its interpretation as a minimax rate is heuristic.

Thus, at present,
\[
R_n^*(\alpha,r)
\le
C n^{-\frac{2(\alpha-1)}{2\alpha+2+d}}
\]
is rigorously known, while the reverse inequality remains open. The conjectured minimax rate is
\[
R_n^*(\alpha,r)
\asymp
n^{-\frac{2(\alpha-1)}{2\alpha+2+d}}
\]
up to logarithmic factors.

\begin{remark}
If future work proves a matching lower bound for the conductivity recovery problem, then Theorem~\ref{thm:state} would upgrade from ``best known upper bound'' to the exact minimax rate.
\end{remark}

\begin{thebibliography}{9}

\bibitem{IOP}
Consistency of Bayesian inference with Gaussian process priors in an elliptic inverse problem.
\href{https://google.iopscience.iop.org/article/10.1088/1361-6420/ab7d2a}{iopscience.iop.org}.

\bibitem{NSS}
Convergence rates for penalized least squares estimators in PDE constrained regression problems.
\href{https://export.arxiv.org/pdf/1809.08818v3.pdf}{arxiv.org}.

\bibitem{MAP}
Convergence rates for the MAP estimator in PDE-regression models with random design.
\href{https://arxiv.org/html/2409.03417}{arxiv.org}.

\bibitem{GLOBAL}
Global polynomial-time estimation in statistical nonlinear inverse problems via generalized stability.
\href{https://arxiv.org/html/2601.09007}{arxiv.org}.

\bibitem{EXA}
Convergence rates for penalized least squares estimators in PDE constrained regression problems.
\href{https://exa.ai/library/publication/rj7jk1yjt53}{exa.ai}.

\end{thebibliography}

\end{document}
```